\section{Application web}\label{sec:application}

\subsection{Architecture}

L'application suit une architecture en couches autour du micro-framework Slim~4 :
\begin{itemize}
	\item \texttt{public/index.php} : point d'entrée unique ;
	\item \texttt{config/} : paramètres (lus depuis \texttt{.env}), conteneur d'injection de dépendances, chaîne de middlewares (erreurs, en-têtes de sécurité, session, CSRF) et routes groupées par rôle ;
	\item \texttt{src/Controller/} : un contrôleur par domaine (stations, locations, abonnements, statistiques\ldots) ;
	\item \texttt{src/Repository/} : toutes les requêtes SQL et les appels de procédures, ce qui isole l'accès aux données ;
	\item \texttt{templates/} : gabarits Twig (mise en page Tailwind, composants réutilisables).
\end{itemize}

La figure~\ref{fig:requete} suit une requête, par exemple l'emprunt d'un vélo. Les middlewares s'appliquent à toutes les routes ; le contrôle de rôle (\texttt{RequireRole}) est attaché à chaque groupe de routes : visiteur, utilisateur connecté, usager, opérateur, administrateur. Le contrôleur ne contient aucune requête SQL : il valide les paramètres, appelle le repository, qui exécute une requête préparée ou une procédure stockée, puis transmet le résultat au gabarit. Une erreur levée par la base (\texttt{SIGNAL}) remonte jusqu'au contrôleur, qui l'affiche comme message à l'utilisateur.

\begin{figure}[htbp]\centering
\begin{tikzpicture}[
	font=\footnotesize,
	etape/.style={draw=neutre, rounded corners=2pt, fill=neutre!6, align=center,
		minimum height=11mm, text width=27mm, inner sep=3pt},
	base/.style={etape, fill=fondvert!35, draw=fondvert!80!black},
	code/.style={font=\scriptsize\ttfamily},
	fleche/.style={-{Stealth[length=2mm]}, thick, neutre},
	node distance=7mm and 6mm,
]
	\node[etape] (nav) {Navigateur\\{\scriptsize\ttfamily POST}\\{\scriptsize\ttfamily /trajets/emprunt}};
	\node[etape, right=of nav] (mw) {Middlewares : erreurs, en-têtes, session, CSRF};
	\node[etape, right=of mw] (route) {Routage et contrôle du rôle (usager)};
	\node[etape, right=of route] (ctrl) {Contrôleur\\{\scriptsize\ttfamily RentalController}};
	\node[etape, below=of ctrl] (repo) {Repository\\{\scriptsize\ttfamily RentalRepository}};
	\node[base, below=of repo] (sp) {MariaDB\\{\scriptsize\ttfamily sp\_emprunter\_velo}\\+ déclencheurs};
	\node[etape, below=of route] (twig) {Redirection et message (succès ou erreur)};
	\node[etape, below=of nav] (rep) {Page affichée\\(gabarit Twig)};
	\draw[fleche] (nav) -- (mw);
	\draw[fleche] (mw) -- (route);
	\draw[fleche] (route) -- (ctrl);
	\draw[fleche] (ctrl) -- (repo);
	\draw[fleche, {Stealth[length=2mm]}-{Stealth[length=2mm]}] (repo) -- node[right, font=\scriptsize, align=left] {appel / résultat\\ou \texttt{SIGNAL}} (sp);
	\draw[fleche] (ctrl.south west) -- (twig.north east);
	\draw[fleche] (twig) -- (rep);
	\draw[fleche] (rep) -- (nav);
\end{tikzpicture}
\caption{Cheminement d'une requête dans l'application}\label{fig:requete}
\end{figure}

\subsection{Fonctionnalités par rôle}

Le tableau~\ref{tab:fonctionnalites} récapitule les pages de l'application et, pour chacune, le traitement réalisé par la base. Les fonctionnalités exigées par le sujet (inscription et connexion, stations en temps réel, gestion des stations, emprunt et restitution, statistiques) sont complétées par la gestion des abonnements et des factures, les signalements, le suivi de la flotte et des tableaux de bord adaptés à chaque rôle.

{\small
\begin{longtable}{>{\raggedright\arraybackslash}p{0.16\linewidth} >{\raggedright\arraybackslash}p{0.38\linewidth} >{\raggedright\arraybackslash}p{0.36\linewidth}}
	\caption{Fonctionnalités de l'application et traitements en base}\label{tab:fonctionnalites}\\
	\toprule
	\textbf{Rôle} & \textbf{Fonctionnalité} & \textbf{Traitement en base} \\
	\midrule\endfirsthead
	\toprule
	\textbf{Rôle} & \textbf{Fonctionnalité} & \textbf{Traitement en base} \\
	\midrule\endhead
	\bottomrule\endlastfoot
	Visiteur & accueil avec la disponibilité du réseau ; inscription ; connexion (tentatives limitées) & vue \texttt{v\_stations\_occupation} ; \texttt{sp\_inscrire\_usager}, \texttt{sp\_authentifier} ; mot de passe bcrypt vérifié par PHP \\
	\midrule
	Usager & stations en temps réel (recherche, filtres) & vue \texttt{v\_stations\_occupation} \\
	& souscription, renouvellement et résiliation d'un abonnement ; factures & \texttt{sp\_souscrire\_abonnement}, \texttt{sp\_renouveler\_abonnement}, \texttt{sp\_resilier\_abonnement} \\
	& emprunt et restitution d'un vélo, avec le coût estimé en direct pendant le trajet & \texttt{sp\_emprunter\_velo}, \texttt{sp\_rendre\_velo} ; estimation par \texttt{fn\_calculer\_cout} \\
	& historique des trajets ; signalement d'un problème ; profil & \texttt{sp\_signaler\_probleme} ; droits limités à certaines colonnes \\
	\midrule
	Opérateur & tableau de bord d'exploitation ; toutes les locations (filtres par date, statut, station, type, usager) & index \texttt{(id\_velo, date\_heure\_emprunt)}, \texttt{(id\_abonnement, date\_heure\_emprunt)} \\
	& flotte : filtres, fiche d'un vélo avec ses historiques d'état et de position & tables \texttt{historique\_etat}, \texttt{historique\_position} \\
	& file des signalements : prise en charge, clôture & index \texttt{(statut, date\_heure)} \\
	& statistiques à onglets, graphiques, export CSV ; propositions de rééquilibrage & six procédures \texttt{sp\_stat\_*} ; \texttt{sp\_stations\_a\_reequilibrer}, \texttt{sp\_proposer\_deplacements} \\
	\midrule
	Administrateur & stations : création, modification, ajout et mise hors service de bornes, suppression & droits d'écriture de \texttt{velo\_app} limités à \texttt{station} et \texttt{borne} \\
	& comptes : recherche, création d'opérateurs & héritage \texttt{utilisateur} / \texttt{operateur} / \texttt{administrateur} \\
\end{longtable}
}

\subsection{Page de statistiques}

La page présente les six statistiques demandées, réparties en trois onglets (\og Trajets \fg{}, \og Stations et régulation \fg{}, \og Recettes et flotte \fg{}) ; seules les procédures de l'onglet affiché sont exécutées. Les résultats sont mis en cache côté application (une heure pour une période passée, cinq minutes si elle inclut le jour courant, trente secondes pour le temps réel), et un lien permet de les recalculer pour mesurer le temps réel des procédures. Chaque rapport est exportable en CSV.
